\section{Diskreetit tekstiesitykset ja klassinen kielimallinnus}\label{sec:discrete}

Kielimallipohjaisessa Text-to-SQL-tehtävässä järjestelmän tavoitteena on muuntaa
luonnollisen kielen tiedontarve symboliseksi SQL-kyselyksi. Tämän prosessin
matemaattinen perusta nojaa kielen esittämiseen laskennallisessa muodossa.

Klassisessa kielimallinnuksessa tekstiä käsitellään diskreetteinä tokeneina. Olkoon
\(\mathcal{V} = \{v_1, v_2, \dots, v_{|\mathcal{V}|}\}\) sanasto eli joukko kaikista mallin
käyttämistä tokeneista ja olkoon
\(\mathcal{W} = (w_1, w_2, \dots, w_{|\mathcal{W}|})\) korpus eli havaittu tokenijono. Tällöin
\(w_t \in \mathcal{V}\) merkitsee korpuksen ajanhetken \(t\) tokenia. Tokenilla voidaan
tarkoittaa esimerkiksi sanaa, osasanaa, merkkiä tai muuta diskreettiä
tekstiyksikköä riippuen käytetystä tokenisointimenetelmästä \cite{jurafsky2023speech}.

Diskreetit tokenit esitetään numeerisessa muodossa, jotta niitä voidaan hyödyntää
tilastollisissa malleissa ja neuroverkoissa. Yksinkertaisin ja yleisesti käytetty
esitystapa on one-hot-enkoodaus. Siinä kukin token \(v_j \in \mathcal{V}\) esitetään
vektorina \(x(v_j) \in \{0,1\}^{1 \times |\mathcal{V}|}\), jonka komponentit ovat
\[
  x(v_j)_k =
  \begin{cases}
    1, & \text{jos } k = j, \\
    0, & \text{muulloin}.
  \end{cases}
\]

Kielimallinnuksessa sanastoa täydennetään usein erikoistokeneilla,
joilla kuvataan sekvenssin rakennetta ja käsitellään harvinaisia havaintoja.
Tällaisia ovat esimerkiksi aloitus- ja lopetusmerkit \(\langle s \rangle\) ja
\(\langle /s \rangle\) sekä tuntemattoman tokenin merkki \(\langle unk \rangle\) \cite{jurafsky2023speech}.

Aloitus- ja lopetusmerkit ovat erityisen hyödyllisiä silloin, kun mallin
ennustettavaa kontekstia halutaan määritellä myös sekvenssin reuna-alueilla.
Esimerkiksi \(n\)-gram-malleissa kontekstin pituus on
vakio, minkä vuoksi sekvenssin alku- ja loppupäähän lisätään tyypillisesti
riittävä määrä erikoistokeneita, jotta kontekstit voidaan muodostaa yhtenäisesti
kaikille ajanhetkille.

Koko korpus voidaan esittää one-hot-vektoreiden jonona
$$
X = (x_1, x_2, \dots, x_{|\mathcal{W}|}),
$$
missä \(x_t = x(w_t) \in \{0,1\}^{1 \times |\mathcal{V}|}\) on tokenin \(w_t\) one-hot-esitys. Jos syöte koostuu
useasta tokenista, voidaan niiden one-hot-esitykset yhdistää esimerkiksi
ketjuttamalla ne yhdeksi vektoriksi
\[
  \tilde{x}_t = [x_{t-n+1} \,\|\, x_{t-n+2} \,\|\, \dots \,\|\, x_{t-1}]
  \in \{0,1\}^{1 \times (n-1)|\mathcal{V}|},
\]
missä \(\|\) merkitsee vektorien vaakasuuntaista ketjuttamista. Tämä muodostaa
luonnollisen syötteen monille kontekstipohjaisille malleille, sillä se säilyttää
kontekstin tokenijärjestyksen eksplisiittisesti.

One-hot-enkoodaus muodostaa perustavanlaatuisen sillan tekstin
ja numeeristen kielimallien välille. Se tarjoaa täsmällisen ja muodollisesti
yksinkertaisen esitystavan, jonka varaan voidaan rakentaa sekä tilastollisia
\(n\)-gram-malleja että myöhempiä neuroverkkopohjaisia kielimalleja.

N-gram-malli on klassinen tilastollinen kielimalli, jossa sekvenssin seuraavan
tokenin todennäköisyyttä estimoidaan rajallisen kontekstin perusteella \cite{jurafsky2023speech}.
Menetelmä perustuu \((n-1)\):n kertaluvun Markov-oletukseen, jonka mukaan
tokenin \(w_t\) todennäköisyys riippuu vain sitä välittömästi edeltävistä
\((n-1)\) tokenista eikä koko aiemmasta historiasta. Markov-oletuksen perusteella
ehdollinen todennäköisyys voidaan approksimoida muodossa
\[
  P(w_t \mid w_1, w_2, \ldots, w_{t-1})
  \approx
  P(w_t \mid w_{t-n+1}, \ldots, w_{t-1}),
\]
missä \(w_t\) on ajanhetken \(t\) tokeni ja
\((w_{t-n+1}, \ldots, w_{t-1})\) muodostaa sitä edeltävän
\((n-1)\) tokenin kontekstin.

N-gram-mallien historialliset juuret yhdistetään usein Shannonin varhaiseen
informaatioteoreettiseen työhön, jossa luonnollista kieltä tarkasteltiin
symbolijonojen tilastollisena prosessina ja eri kertalukujen approksimaatioita
havainnollistettiin englannin kielellä
\cite{shannon1948mathematical,shannon1951prediction}.

Olkoon \(\mathcal{V} = \{v_1, v_2, \ldots, v_{|\mathcal{V}|}\}\) sanasto, jossa \(|\mathcal{V}|\) on sanaston
koko, ja \(\mathcal{W} = (w_1, w_2, \ldots, w_{|\mathcal{W}|})\) havaintoaineisto. N-gram-malli
voidaan toteuttaa muodostamalla lukumatriisi
\(C \in \mathbb{R}^{|\mathcal{V}|^{n-1} \times |\mathcal{V}|}\), jossa \(C_{c,j}\) edustaa sitä,
kuinka monta kertaa kontekstia \(c\) seuraa tokeni \(v_j\)
havaintoaineistossa. Formaalisti
\[
  C_{c,j}
  =
  \left|
  \left\{
    t \in \mathbb{N}
    \mid
    (w_{t-n+1}, \ldots, w_{t-1}) = c
    \ \text{ja}\
    w_t = v_j
  \right\}
  \right|.
\]

Kun lukumatriisin \(C\) rivit normalisoidaan jakamalla jokainen alkio rivinsä
summalla, saadaan frekvenssimatriisi
\(F \in \mathbb{R}^{|\mathcal{V}|^{n-1} \times |\mathcal{V}|}\), jossa
\[
  F_{c,j}
  =
  \frac{C_{c,j}}{\sum_{k=1}^{|\mathcal{V}|} C_{c,k}},
\]
ja joka edustaa ehdollista todennäköisyyttä
\[
  P(w_t = v_j \mid (w_{t-n+1}, \ldots, w_{t-1}) = c).
\]

Tekstin generointi n-gram-mallilla tapahtuu iteratiivisesti hyödyntämällä
frekvenssimatriisia \(F\). Ensin valitaan alkukonteksti, esimerkiksi
erikoismerkeillä aloitettu \((n-1)\)-tokenin pituinen alkujakso. Tämän jälkeen
mallin tämänhetkistä kontekstia vastaava rivi \(F_{c,\cdot}\) tulkitaan
todennäköisyysjakaumaksi, josta seuraava tokeni valitaan satunnaisotannalla.
Kun tokeni \(w_t\) on generoitu, konteksti päivitetään siirtämällä ikkuna yhden
askeleen eteenpäin siten, että uusi konteksti muodostuu edeltävistä
\((n-1)\) generoidusta tokenista. Prosessia jatketaan, kunnes generoidaan
päätemerkki tai saavutetaan haluttu tekstin pituus.

Diskreetit esitykset ja frekvenssipohjaiset \(n\)-gram-mallit tarjoavat
käsitteellisesti selkeän lähtökohdan kielimallinnukselle, mutta niihin liittyy
kaksi perustavanlaatuista rajoitetta.

Ensinäkin one-hot-enkoodaus ilmaisee vain tokenin identiteetin. Koska erisuurten tokeneiden vektorit ovat keskenään ortogonaalisia
\[
  v_i \neq v_j \implies x(v_i) x(v_j)^T = 0,
\]
niiden kosinisimilariteetti on poikkeuksetta nolla ja euklidinen etäisyys vakio \(\sqrt{2}\). Tästä johtuen esitys ei voi mallintaa tokeneiden välisiä semanttisia tai syntaktisia suhteita \cite{mikolov2013efficient}.

Toiseksi kontekstikoon \(n\) kasvaessa lukumatriisista
\(C \in \mathbb{R}^{|\mathcal{V}|^{n-1} \times |\mathcal{V}|}\) tulee nopeasti erittäin suuri ja
käytännöllisissä sovelluksissa hyvin harva, koska vain pieni osa kaikista
mahdollisista konteksteista esiintyy tyypillisessä havaintoaineistossa. Tämä
lisää sekä muistin- että laskentatehon tarvetta
\cite{jurafsky2023speech}.

Tämän vuoksi on hyödyllistä tarkastella samaa ongelmaa
optimointitehtävänä, mikä muodostaa luontevan siirtymän kohti jatkuvia
vektoriesityksiä ja parametrisoituja neuroverkkopohjaisia malleja.
